% GENERATED FILE. Edit canonical lessons, then run npm run generate:books.
% canonical-source-hash: fnv1a64:d3b2a7c286ca97ed
% canonical-lessons: RU-C294-vinovnyy, RU-C294-nakazaniye, RU-C294-nakazyvat, RU-C294-priznavatsya, RU-C294-opravdyvat, RU-R294-review-287-294

\chapter{Guilt and punishment}
\label{ch:ru-guilt-and-punishment}
\hlchaptermodality{294}

\begin{chapteropening}
\textbf{By the end of this chapter:} \emph{I can say that someone was guilty, was punished, confessed or was acquitted.}

Complete the last lesson of chapter 294: I can say that someone was guilty, was punished, confessed or was acquitted.
\end{chapteropening}

\section[\texorpdfstring{\ru{виновный} (vinóvnyy)}{виновный (vinóvnyy)}]{\ru{виновный} (vinóvnyy) — guilty}

\label{lesson:RU-C294-vinovnyy}



\begin{quote}
Before the new one: say the Russian for a prison, then the Russian for to accuse.
\end{quote}

\subsection*{You'll want to know: \ru{виновный}}
\textbf{\ru{виновный}} (\emph{vinóvnyy}) — \textquotedblleft{}guilty\textquotedblright{}.

\begin{quote}
\ru{Суд} \ru{признал} \ru{его} \ru{виновным}.
\end{quote}

\begin{quote}
\emph{Sud priznál yevó vinóvnym.}
\end{quote}

\begin{quote}
The court found him guilty.
\end{quote}

\emph{\ru{Виновный} \ru{в}} + prepositional. In speech, of a small fault: \emph{\ru{я} \ru{виноват}} — it is my fault.

\begin{cousinweb}
\emph{\ru{виновный}} is built on \emph{\ru{вина}} \textquotedblleft{}guilt\textquotedblright{}, an inherited Slavic word, with the adjective ending \emph{-\ru{ный}}; EtymDB lists no parent.
\end{cousinweb}

\begin{grammarlens}[title={What you've built}]
One more word for telling, explaining and describing at B1.
\end{grammarlens}

\subsection*{Your turn}
\begin{itemize}
\raggedright
  \item \emph{Say it:} \emph{vinóvnyy}
  \item \emph{Say it:} \emph{vinóvnyy}, once more
  \item \emph{Say it:} a sentence of your own with \emph{vinóvnyy}
  \item \emph{Recall:} say the Russian for a prison, then the Russian for to accuse, then say \emph{vinóvnyy} again
\end{itemize}

\begin{tcolorbox}[breakable,colback=teal!4,colframe=teal!35!black,title={Before you move on}]
What does \ru{виновный} mean? (\textquotedblleft{}Guilty\textquotedblright{}.) Say it once more.
\end{tcolorbox}
\section[\texorpdfstring{\ru{наказание} (nakazániye)}{наказание (nakazániye)}]{\ru{наказание} (nakazániye) — a punishment}

\label{lesson:RU-C294-nakazaniye}



\begin{quote}
Before the new one: say the Russian for to accuse, then the Russian for guilty.
\end{quote}

\subsection*{You'll want to know: \ru{наказание}}
\textbf{\ru{наказание}} (\emph{nakazániye}) — \textquotedblleft{}a punishment\textquotedblright{}.

\begin{quote}
\ru{Наказание} \ru{было} \ru{справедливым}.
\end{quote}

\begin{quote}
\emph{Nakazániye býlo spravedlívym.}
\end{quote}

\begin{quote}
The punishment was fair.
\end{quote}

Neuter.

\begin{cousinweb}
\emph{\ru{наказание}} is built on \emph{\ru{наказать}}, from \emph{\ru{на}-} and \emph{\ru{казать}} \textquotedblleft{}to show, to tell\textquotedblright{}, inherited from Proto-Slavic \emph{kazati} (EtymDB): first \textquotedblleft{}an instruction\textquotedblright{}, then a punishment.
\end{cousinweb}

\begin{grammarlens}[title={What you've built}]
One more word for telling, explaining and describing at B1.
\end{grammarlens}

\subsection*{Your turn}
\begin{itemize}
\raggedright
  \item \emph{Say it:} \emph{nakazániye}
  \item \emph{Say it:} \emph{nakazániye}, once more
  \item \emph{Say it:} a sentence of your own with \emph{nakazániye}
  \item \emph{Recall:} say the Russian for to accuse, then the Russian for guilty, then say \emph{nakazániye} again
\end{itemize}

\begin{tcolorbox}[breakable,colback=teal!4,colframe=teal!35!black,title={Before you move on}]
What does \ru{наказание} mean? (\textquotedblleft{}A punishment\textquotedblright{}.) Say it once more.
\end{tcolorbox}
\section[\texorpdfstring{\ru{наказывать} (nakázyvat)}{наказывать (nakázyvat)}]{\ru{наказывать} (nakázyvat) — to punish}

\label{lesson:RU-C294-nakazyvat}



\begin{quote}
Before the new one: say the Russian for guilty, then the Russian for a punishment.
\end{quote}

\subsection*{You'll want to know: \ru{наказывать}}
\textbf{\ru{наказывать}} (\emph{nakázyvat}) — \textquotedblleft{}to punish\textquotedblright{}.

\begin{quote}
\ru{Родители} \ru{редко} \ru{наказывали} \ru{сына}.
\end{quote}

\begin{quote}
\emph{Rodíteli rédko nakázyvali sýna.}
\end{quote}

\begin{quote}
The parents rarely punished their son.
\end{quote}

Imperfective; perfective \emph{\ru{наказать}} (\emph{nakazát}), \emph{\ru{накажу}}, \emph{\ru{накажешь}}. The reason takes \emph{\ru{за}} + accusative: \emph{\ru{наказать} \ru{за} \ru{ложь}}.

\begin{cousinweb}
\emph{\ru{наказывать}} is built on \emph{\ru{на}-} and \emph{\ru{казать}} \textquotedblleft{}to show, to tell\textquotedblright{}, inherited from Proto-Slavic \emph{kazati} (EtymDB): first \textquotedblleft{}to instruct\textquotedblright{}.
\end{cousinweb}

\begin{grammarlens}[title={What you've built}]
One more word for telling, explaining and describing at B1.
\end{grammarlens}

\subsection*{Your turn}
\begin{itemize}
\raggedright
  \item \emph{Say it:} \emph{nakázyvat}
  \item \emph{Say it:} \emph{nakázyvat}, once more
  \item \emph{Say it:} a sentence of your own with \emph{nakázyvat}
  \item \emph{Recall:} say the Russian for guilty, then the Russian for a punishment, then say \emph{nakázyvat} again
\end{itemize}

\begin{tcolorbox}[breakable,colback=teal!4,colframe=teal!35!black,title={Before you move on}]
What does \ru{наказывать} mean? (\textquotedblleft{}To punish\textquotedblright{}.) Say it once more.
\end{tcolorbox}
\section[\texorpdfstring{\ru{признаваться} (priznavátsya)}{признаваться (priznavátsya)}]{\ru{признаваться} (priznavátsya) — to confess, to admit}

\label{lesson:RU-C294-priznavatsya}



\begin{quote}
Before the new one: say the Russian for a punishment, then the Russian for to punish.
\end{quote}

\subsection*{You'll want to know: \ru{признаваться}}
\textbf{\ru{признаваться}} (\emph{priznavátsya}) — \textquotedblleft{}to confess, to admit\textquotedblright{}.

\begin{quote}
\ru{Он} \ru{признался}, \ru{что} \ru{взял} \ru{деньги}.
\end{quote}

\begin{quote}
\emph{On priználsya, shto vzyal déngi.}
\end{quote}

\begin{quote}
He admitted that he had taken the money.
\end{quote}

Imperfective; perfective \emph{\ru{признаться}} (\emph{priznátsya}). Present \emph{\ru{признаюсь}}, \emph{\ru{признаёшься}} (\emph{priznayóshsya}). \emph{\ru{Признаться} \ru{в} \ru{любви}} — to declare one's love.

\begin{cousinweb}
\emph{\ru{признаваться}} is built on \emph{\ru{при}-} and \emph{\ru{знать}} \textquotedblleft{}to know\textquotedblright{}, which is inherited from Proto-Slavic \emph{znati}, from PIE \emph{ǵneh\textsubscript{3}-} \textquotedblleft{}to know\textquotedblright{} (EtymDB): to make known.
\end{cousinweb}

\begin{grammarlens}[title={What you've built}]
One more word for telling, explaining and describing at B1.
\end{grammarlens}

\subsection*{Your turn}
\begin{itemize}
\raggedright
  \item \emph{Say it:} \emph{priznavátsya}
  \item \emph{Say it:} \emph{priznavátsya}, once more
  \item \emph{Say it:} a sentence of your own with \emph{priznavátsya}
  \item \emph{Recall:} say the Russian for a punishment, then the Russian for to punish, then say \emph{priznavátsya} again
\end{itemize}

\begin{tcolorbox}[breakable,colback=teal!4,colframe=teal!35!black,title={Before you move on}]
What does \ru{признаваться} mean? (\textquotedblleft{}To confess, to admit\textquotedblright{}.) Say it once more.
\end{tcolorbox}
\section[\texorpdfstring{\ru{оправдывать} (oprávdyvat)}{оправдывать (oprávdyvat)}]{\ru{оправдывать} (oprávdyvat) — to acquit; to justify}

\label{lesson:RU-C294-opravdyvat}



\begin{quote}
Before the new one: say the Russian for to punish, then the Russian for to confess.
\end{quote}

\subsection*{You'll want to know: \ru{оправдывать}}
\textbf{\ru{оправдывать}} (\emph{oprávdyvat}) — \textquotedblleft{}to acquit; to justify\textquotedblright{}.

\begin{quote}
\ru{Суд} \ru{оправдал} \ru{молодого} \ru{человека}.
\end{quote}

\begin{quote}
\emph{Sud opravdál molodóvo chelovéka.}
\end{quote}

\begin{quote}
The court acquitted the young man.
\end{quote}

Imperfective; the perfective is \emph{\ru{оправдать}} (\emph{opravdát}). Also \emph{\ru{оправдать} \ru{надежды}} — to live up to hopes.

\begin{cousinweb}
\emph{\ru{оправдывать}} is built on \emph{\ru{о}-} and \emph{\ru{правда}} \textquotedblleft{}truth\textquotedblright{}, from \emph{\ru{правый}} \textquotedblleft{}right, fair\textquotedblright{}, an inherited Slavic word: to show someone is in the right. EtymDB lists no parent.
\end{cousinweb}

\begin{grammarlens}[title={What you've built}]
One more word for telling, explaining and describing at B1.
\end{grammarlens}

\subsection*{Your turn}
\begin{itemize}
\raggedright
  \item \emph{Say it:} \emph{oprávdyvat}
  \item \emph{Say it:} \emph{oprávdyvat}, once more
  \item \emph{Say it:} a sentence of your own with \emph{oprávdyvat}
  \item \emph{Recall:} say the Russian for to punish, then the Russian for to confess, then say \emph{oprávdyvat} again
\end{itemize}

\begin{tcolorbox}[breakable,colback=teal!4,colframe=teal!35!black,title={Before you move on}]
What does \ru{оправдывать} mean? (\textquotedblleft{}To acquit; to justify\textquotedblright{}.) Say it once more.
\end{tcolorbox}
\section[(dialogue)]{Words from chapters 287–294 — a review}

\label{lesson:RU-R294-review-287-294}



\begin{quote}
Say the Russian for a crime and a criminal.
\end{quote}

\subsection*{Your turn}
Say these aloud:

\begin{itemize}
\raggedright
  \item a crime, a criminal, a victim, a clue, a witness
  \item a theft, a thief, to steal, a pickpocket, a break-in
  \item a robbery, a robber, to rob, an attack, to attack
  \item a swindler, to deceive, to hijack, to kidnap, to hide
  \item an investigator, an investigation, to investigate, a suspect, to suspect
  \item to search, to arrest, to question, questioning, a statement
  \item a court, a lawyer, a sentence, a prison, to accuse
  \item guilty, a punishment, to punish, to confess, to acquit
\end{itemize}

\begin{tcolorbox}[breakable,colback=teal!4,colframe=teal!35!black,title={Before you move on}]
A crime? (\textbf{\ru{преступление}}.) A criminal? (\textbf{\ru{преступник}}.) A victim? (\textbf{\ru{жертва}}.) A clue? (\textbf{\ru{улика}}.) A witness? (\textbf{\ru{свидетель}}.) A theft? (\textbf{\ru{кража}}.) A thief? (\textbf{\ru{вор}}.) To steal? (\textbf{\ru{красть}}.) A pickpocket? (\textbf{\ru{карманник}}.) A break-in? (\textbf{\ru{взлом}}.) A robbery? (\textbf{\ru{ограбление}}.) A robber? (\textbf{\ru{грабитель}}.) To rob? (\textbf{\ru{грабить}}.) An attack? (\textbf{\ru{нападение}}.) To attack? (\textbf{\ru{нападать}}.) A swindler? (\textbf{\ru{мошенник}}.) To deceive? (\textbf{\ru{обманывать}}.) To hijack? (\textbf{\ru{угонять}}.) To kidnap? (\textbf{\ru{похищать}}.) To hide? (\textbf{\ru{прятаться}}.) An investigator? (\textbf{\ru{следователь}}.) An investigation? (\textbf{\ru{расследование}}.) To investigate? (\textbf{\ru{расследовать}}.) A suspect? (\textbf{\ru{подозреваемый}}.) To suspect? (\textbf{\ru{подозревать}}.) To search? (\textbf{\ru{обыскивать}}.) To arrest? (\textbf{\ru{арестовывать}}.) To question? (\textbf{\ru{допрашивать}}.) Questioning? (\textbf{\ru{допрос}}.) A statement? (\textbf{\ru{показания}}.) A court? (\textbf{\ru{суд}}.) A lawyer? (\textbf{\ru{адвокат}}.) A sentence? (\textbf{\ru{приговор}}.) A prison? (\textbf{\ru{тюрьма}}.) To accuse? (\textbf{\ru{обвинять}}.) Guilty? (\textbf{\ru{виновный}}.) A punishment? (\textbf{\ru{наказание}}.) To punish? (\textbf{\ru{наказывать}}.) To confess? (\textbf{\ru{признаваться}}.) To acquit? (\textbf{\ru{оправдывать}}.)
\end{tcolorbox}
